\documentclass[10pt]{article}
\usepackage[margin=0.85in]{geometry}
\usepackage{amsmath,amssymb,amsthm}
\usepackage{hyperref}
\newtheorem{theorem}{Theorem}
\newtheorem{lemma}{Lemma}
\title{A deterministic boundary counterexample to conjecture 00000003479}
\author{earthking11}
\date{October 6, 2026}
\begin{document}
\maketitle
\section*{Statement and interpretation}
The conjecture defines a harmonious coloring to be a proper coloring in
which sums of endpoint color labels distinguish all unordered edges. It
claims that the harmonious chromatic number $h(G(n,p))$ concentrates at
$(1+o(1))\sqrt{np}$, uniformly for $p\geq \log n/n$. We use precisely the
sum definition, not the different standard definition based on unordered
color pairs. Colors in a $k$-coloring have natural labels $0,\ldots,k-1$.
The argument uses only properness, so any additional sum restriction
cannot rescue the claimed upper concentration.

\begin{lemma}
Every harmonious $k$-coloring of $K_n$ has $k\geq n$.
\end{lemma}
\begin{proof}
Every two distinct vertices of $K_n$ are adjacent. Properness therefore
makes the color map injective. An injection from $n$ vertices into $k$
colors implies $n\leq k$.
\end{proof}

\begin{theorem}
For every integer $n>4$, the probability that $G(n,1)$ admits a harmonious
$k$-coloring with $k\leq 2\sqrt n$ is zero. Consequently, the stated
constant-one concentration law is false.
\end{theorem}
\begin{proof}
At $p=1$, every edge is present with probability one and $G(n,1)=K_n$
almost surely. This parameter is admissible for every $n\geq1$, because
$\log n\leq n-1\leq n$.
By the lemma, a coloring in the asserted event would give
$n\leq k\leq2\sqrt n$. Squaring implies $n^2\leq4n$, which contradicts
$n>4$. Thus the upper event has probability exactly zero for all $n>4$.

In particular, if $h$ is the minimum number of colors, the relative
concentration event
\[
 \left|\frac{h(G(n,1))}{\sqrt n}-1\right|<1
\]
is contained in that impossible upper event. Its probability is eventually
zero and tends to zero, not one. The conjectured relative concentration
would require convergence to one for every positive tolerance, including
this tolerance. Failure at the admissible deterministic boundary refutes
the claimed uniform law. No finite-size numerical anchor is used.
\end{proof}

\section*{Formal verification and scope}
The Lean project uses Mathlib's actual \texttt{SimpleGraph}, graph
\texttt{Coloring}, and probability mass functions. It defines harmonious
colorings with the endpoint-sum condition. For every finite graph it
constructs a coloring with labels $(2n^2+1)i+i^2$: a label-pair sum determines
the index sum and square sum, hence its unordered endpoint pair. Thus a
minimum exists. Lean defines the minimum by natural-number well-ordering,
proves its attainment and minimality, and identifies its relative-error
event with the formally defined concentration event.
The definition \texttt{GnpOneLaw} specifies that each distinct edge has
probability one; Lean proves every graph in its support is $K_n$.
The explicit point mass at $K_n$ satisfies this law. Thus the proof does
not assume a counting bound or a random-graph tail estimate.
The capstone \texttt{conjecture\_boundary\_is\_false} proves the negation
of constant-one relative concentration at $p=1$. All statements compile
on Lean 4.33.1 with pinned Mathlib, without unfinished proofs, native
evaluation, or additional axioms. The only reported foundational axioms
are propositional extensionality, classical choice, and quotient soundness.
We do not claim a corrected concentration law for $0<p<1$, nor need the
conjecture's additional low-probability clause.
\end{document}
